\documentclass[manuscript, screen, review]{acmart} % Use [acmsmall, screen] when ready.

\usepackage{acmpaper}

% For the Cahier du GERAD.
%\renewcommand{\year}{2020}     % If you don't want the current year.
\newcommand{\cahiernumber}{00}  % Insert your Cahier du GERAD number.

%\setcopyright{acmlicensed}
%\copyrightyear{2018}
%\acmYear{2018}
%\acmDOI{XXXXXXX.XXXXXXX}

% Insert the name of "your journal" with
\acmJournal{TOMS}

\begin{document}

\title{ExactPenalty: an Exact Penalty approach for Large-Scale Nonlinear Optimization}

\author{Youssef Diouane}
\affiliation{%
  \institution{GERAD and Polytechnique Montr\'eal},
  \country{Canada}
}
\email{youssef.diouane@polymtl.ca}

\author{Maxence Gollier}
\affiliation{%
  \institution{GERAD and Polytechnique Montr\'eal},
  \country{Canada}
}
\email{maxence-2.gollier@polymtl.ca}

\author{Dominique Orban}
\affiliation{%
  \institution{GERAD and Polytechnique Montr\'eal},
  \country{Canada}
}
\email{dominique.orban@gerad.ca}


\begin{abstract}
  Insert abstract here. 
\end{abstract}

\keywords{Keyword1, Keyword2, Keyword3}

%\received{20 February 2007}
%\received[revised]{12 March 2009}
%\received[accepted]{5 June 2009}

%%
%% This command processes the author and affiliation and title
%% information and builds the first part of the formatted document.
\maketitle

\section{Introduction}

This is the introduction. An optimization problem has the general form
\begin{equation}
  \label{eq:opt-problem}
  \minimize{x \in \R^n} \ f(x) \quad \st \ \ell \leq c(x) \leq u.
\end{equation}

For the real stuff, check \cite{wright-orban-2002}.

\section{Conclusion}

Concusions will come in due time.
\smarttodo{Come up with witty conclusions}

\begin{acks}
To Robert, for the bagels and explaining CMYK and color spaces.
\end{acks}

%%
%% The next two lines define the bibliography style to be used, and
%% the bibliography file.
\bibliographystyle{abbrvnat}
\bibliography{abbrv, \jobname}

\end{document}